% ============================================================
% Related work moved from chapter infonomy (chapters/infonomy.tex,
% "Extrapolating volition with recursive information markets"),
% verbatim from the paper's \subsection{Related work}.
% Uses chapter-local macros \citeps, \citets (aliases of \citep/\citet,
% defined in macros/infonomy.tex); \bleg is global (preamble/thesis.tex).
% ============================================================

\subsection{Related work}
\label{infonomy:sec:related}

% The information bazaar was introduced by \citets{weissRedesigningInformationMarkets2024}. Apart from this, several other streams of research are tangent to our work.

\paragraph{Information economics.} The foundational theory of information economics\bleg{feel like this might elicit some annoying comments from reviewers} is the \emph{value-of information} framework \citeps{kuhn1953information, howard1966information, raiffa1961AppliedStatisticalDecisionTheory}, which treats information as an instrumental good \citeps{vanalstyneProposalValuingInformation1999}. The buyer's inspection paradox was introduced by \citets{arrow1972economic} and named by \citets{vanalstyneProposalValuingInformation1999}. \citets{hansonIPBarbedWire2011,hansonRahEfficientIP2011} discussed the inefficiency of information markets in the context of intellectual property (IP) law, commenting: ``\emph{just as farmers developed barbed-wire, someday I expect IP advocates will develop better forms of intellectual property}''.

\paragraph{Mechanism design for information markets.} Simple, naive information markets suffer from a number of flaws: \emph{the low cost of duplication} \citeps{samuelson2009economics}, \emph{the transaction costs of tenders} \citeps{tenders}, the transaction cost of \emph{learning} new information, and \emph{information asymmetry} \citeps{arrow1972economic,vanalstyneProposalValuingInformation1999}. \citets{conitzerPredictionMarketsMechanism2012} introduced a mechanism for rewarding information-providing agents based on their influence on prediction market prices, though this is only applicable in contexts where ground-truth (prediction market resolutions) is available.


\paragraph{Scalable oversight} The fundamental limitation of RLHF that it relies on a human's ability to judge a (potentially superhuman) AI's outputs has long been recognized \citeps{christiano2018IDA}, and is known as the \emph{scalable oversight} problem in the AI alignment literature \citeps{hubinger2020alignmentProposalsComplexityClasses,bowmanMeasuringProgressScalable2022}. One well-known scalable oversight proposal is Debate \citeps{irvingAISafetyDebate2018}; our proposal to augment RLHF with information markets may be seen as yet another such proposal.

\paragraph{Mechanism design with LLMs.} Market design for LLM participants has opened up many new frontiers previously not possible with humans alone. Apart from the information bazaar \citeps{weissRedesigningInformationMarkets2024}, this includes e.g. token auctions for online advertising \citeps{duttingMechanismDesignLarge2024}, economic simulations with AI agents \citeps{liEconAgentLargeLanguage2024a}, and forecasting with LLMs \citeps{schoeneggerWisdomSiliconCrowd2024,halawiApproachingHumanLevelForecasting2024,palekaConsistencyChecksLanguage2024,buterin2024InfoFinance}.

\paragraph{Zero-knowledge proofs.} Another way for a seller to prove the value of his information without revealing it is via a \emph{zero-knowledge proof} \citeps{goldreich2001FoundationsCryptographyV1} -- in formal settings, this is available if the information is a solution to a PSPACE problem \citeps{zkPSPACE1,zkPSPACE2}; extending this to informal settings is an active area of work \citeps{hammondNeuralInteractiveProofs2024}.

\paragraph{Miscellaneous} Recent works like \citets{manelli2006Bundling,babiaoff2012OptimalMechanismsSellingInformation,bergemannDesignPriceInformation2018} have studied ``optimal mechanisms for selling information'', but from the point-of-view of a seller maximizing his revenue -- we, on the other hand, are interested in improving the efficiency of the information market itself. Information market mechanisms have also been designed for \emph{data markets} in machine learning, e.g. \citets{fallah2024ThreeLayerDataMarkets,chenSellingDataMachine2022,ghorbani2019DataShapley}.\bleg{all papers in this para are totally irrelevant but I suspect reviewers may bring them up}
